\chapter{How Machine Learning Works}
\label{ch:learning}

\chapterguide%
  {\chapref{ch:data}: a trustworthy table, a clear observation unit, and a split that matches the real prediction setting.}
  {distinguish learning signals, fitting objectives, evaluation, and the mathematics each workflow needs.}

The available feedback determines the learning objective and the mathematics needed.

\begin{keyidea}
Labels guide prediction; unlabeled data supports structure discovery; rewards guide actions.
\end{keyidea}

\section{Start with the learning signal}

\begin{mlfigure}
\begin{tikzpicture}[
  q/.style={draw=MLNavy,fill=MLNavy,text=white,rounded corners=2mm,text width=3.2cm,minimum height=0.9cm,align=center,font=\bfseries\small},
  bx/.style={draw=MLBlue,fill=MLPaleBlue,rounded corners=2mm,text width=3.0cm,minimum height=1.0cm,align=center,font=\small\color{MLNavy}},
  leaf/.style={draw=MLTeal,fill=MLPaleTeal,rounded corners=2mm,text width=2.8cm,minimum height=1.2cm,align=center,font=\scriptsize\color{MLNavy}},
  arr/.style={-{Stealth[length=2mm]},thick,draw=MLMuted}
]
\node[q] (s) at (0,3.6) {What feedback do we have?};
\node[bx] (full) at (-5.1,1.8) {Labels for examples};
\node[bx] (few) at (-1.7,1.8) {Only some labels};
\node[bx] (none) at (1.7,1.8) {No external label};
\node[bx] (reward) at (5.1,1.8) {Rewards after actions};
\node[leaf] (sup) at (-5.1,-0.2) {\textbf{Supervised}\\learn $X\rightarrow y$};
\node[leaf] (semi) at (-1.7,-0.2) {\textbf{Limited-label}\\combine labeled and unlabeled data};
\node[leaf] (unsup) at (1.7,-0.2) {\textbf{Unsupervised}\\discover structure in $X$};
\node[leaf] (rl) at (5.1,-0.2) {\textbf{Reinforcement}\\learn behavior from interaction};
\draw[arr] (s.south) -- (0,2.75); \draw[draw=MLMuted,thick] (-5.1,2.75)--(5.1,2.75);
\foreach \x in {-5.1,-1.7,1.7,5.1} {\draw[arr] (\x,2.75)--(\x,2.3);}
\draw[arr] (full)--(sup); \draw[arr] (few)--(semi); \draw[arr] (none)--(unsup); \draw[arr] (reward)--(rl);
\end{tikzpicture}
\end{mlfigure}

\begin{mltable}
\begin{tabularx}{\linewidth}{@{}>{\bfseries\leavevmode\color{MLNavy}}p{0.21\linewidth}p{0.22\linewidth}X X@{}}
\toprule
\rowcolor{MLPaleBlue!92}
Setting & Given & What is learned? & Typical question \\
\midrule
Supervised & $X$ and label $y$ & A mapping from features to a label & ``What value or class should I predict?'' \\
Unsupervised & $X$ only & Structure, representation, density, or unusualness in the features & ``What patterns are present when no answer key is supplied?'' \\
Semi-supervised / active & $X$ plus some labels & A predictive mapping that also uses unlabeled structure or selectively requests labels & ``How can scarce labels be used efficiently?'' \\
Reinforcement & States, actions, rewards, transitions & A policy or value rule for sequential decisions & ``Which action leads to good long-term outcomes?'' \\
\bottomrule
\end{tabularx}
\end{mltable}

\section{The common skeleton}

\begin{mlfigure}
\begin{tikzpicture}[
  node distance=0.3cm,
  bx/.style={draw=MLBlue,fill=MLPaleBlue,rounded corners=1.5mm,minimum width=2.65cm,minimum height=0.76cm,inner xsep=3pt,align=center,font=\scriptsize\bfseries\color{MLNavy}},
  arr/.style={-{Stealth[length=2mm]},thick,draw=MLTeal}
]
\node[bx] (a) {Define the question};
\node[bx,right=of a] (b) {Prepare the data};
\node[bx,right=of b] (c) {Fit the learning rule};
\node[bx,right=of c] (d) {Evaluate the result};
\node[bx,right=of d] (e) {Use and monitor};
\draw[arr] (a)--(b); \draw[arr] (b)--(c); \draw[arr] (c)--(d); \draw[arr] (d)--(e);
\end{tikzpicture}
\end{mlfigure}

Supervised learning compares predictions with labels. Unsupervised learning needs evidence beyond an answer key. Reinforcement learning evaluates behavior across trajectories.

\section{What changes between supervised and unsupervised learning}

\begin{mltable}
\begin{tabularx}{\linewidth}{@{}>{\bfseries\leavevmode\color{MLNavy}}p{0.25\linewidth}X X@{}}
\toprule
\rowcolor{MLPaleBlue!92}
Question & Supervised learning & Unsupervised learning \\
\midrule
Data? & Features $X$ and labels $y$ & Features $X$ only \\
Fitting objective? & Minimize prediction loss & Cluster compactness, reconstruction, variance, density, or isolation \\
Output? & Predictions, scores, or probabilities & Groups, representations, densities, anomalies, or patterns \\
Baseline? & Mean, median, majority, or domain rule & A trivial structure to beat: one cluster, a random partition, or the mean as the reconstruction \\
What counts as ``good''? & Strong performance on realistic unseen data & Structure that is stable, makes sense to a domain expert, and is useful for a downstream task \\
\bottomrule
\end{tabularx}
\end{mltable}

\begin{pitfall}
A cluster need not be a class, a PCA component is not a prediction, and an anomaly score is not a fraud probability. Choose evaluation to match the objective.
\end{pitfall}

\section{A first look at optimization: improve the objective step by step}
\label{sec:gradient-descent}

Some models fit directly; others iteratively reduce an objective $J(\bm{\theta})$. Gradient descent updates parameters by
\[
  \bm{\theta}\leftarrow\bm{\theta}-\alpha\nabla J(\bm{\theta}),
\]
where $\nabla J$ points in the direction of steepest increase and the learning rate $\alpha$ controls the step size.

\begin{intuition}
The gradient sets direction; the learning rate sets stride. Small steps may be slow; large steps may overshoot.
\end{intuition}

\begin{workedexample}
For $J(w)=(w-4)^2$, $J'(w)=2(w-4)$ and the minimum is $w=4$. Starting at $w=0$ with $\alpha=0.1$ gives $w\leftarrow0-0.1(-8)=0.8$. Repeated updates approach $4$; \chapref{ch:foundations} explains the gradient geometry.
\end{workedexample}

\begin{keyidea}
Reducing training loss does not establish generalization; evaluate on unseen data.
\end{keyidea}

\section{Now learn only the mathematics that the workflow needs}

Start with center, spread, probability, vectors, matrix shapes, dot products, distances, scaling, and gradients. Revisit deeper mathematics when needed.

\begin{quickcheck}
\begin{enumerate}
  \item What feedback signal separates supervised, unsupervised, limited-label, and reinforcement-learning problems?
  \item Why can two learning settings share data preparation and monitoring practices while requiring different fitting objectives?
  \item What is the difference between reducing a training objective and demonstrating generalization?
\end{enumerate}
\end{quickcheck}
